\begin{definition}
\label{Defn:NearlyVertexTra}
The graph $(V, \widetilde E)$ is said to be \textnormal{nearly vertex-transitive under a group}~$\calG$ if $\calG$ acts transitively on $V$ from the left, no index two subgroup of $\calG$ acts transitively on $V$, the set 
$
\widetilde \calN (\widetilde \calN(g(\scrA)))$ 
is contained in 
$g(\widetilde \calN(\widetilde \calN(\scrA)))
$ 
for any $g\in \calG$, and for each~$\theta_i, 1\leq i\leq d$ and $v\in V$, there is an automorphism or an anti-automorphism $\psi_{i,v}$ of the group $\calG$ such that one of 
$$\theta_i(g\cdot v) =  \psi_{i,v}(g) \cdot \theta_i(v), 
\quad 
\theta_i(g\cdot v) =  \psi_{i, v}(g^\mo) \cdot \theta_i(v)$$
holds for any $g\in \calG$. In this situation, we say that $(V, \widetilde E)$ is \textnormal{nearly vertex-transitive under $\calG$ with respect to} $\{\psi_{i,v}\}$.
\end{definition}
